\chapter{Huh?}

\section{P, NP}

La classe $P$ è la classe dei linguaggi riconoscibili in tempo deterministico polinomiale

$$
P = \bigcup
$$
%
La classe $NP$ è la classe dei linguaggi riconoscibili in tempo nondeterministico polinomiale

$$
NP = \bigcup
$$

\subsection{Teorema di inclusione deterministica: $P \subseteq NP$}
$$
P \subseteq NP,
$$
inoltre, se $\mathcal{L} \in NP$, allora $\mathcal L$ è decidibile.

\section{Sull'output}

\subsection{Lunghezza dell'output}

La lunghezza dell'output è sempre minore o uguale alla lunghezza dell'input + lunghezza della computazione + 1.

\subsection{Output prodotto in tempo \textit{n}}
\paragraph{Dimostrazione}
Supponiamo $(X,Y) \in \sigma(\mathfrak{M})$

\subsection{Output prodotto in tempo \textit{t}}

\section{Chiusura per composizione}

Siano $f$ e $g$ funzioni calcolabili in tempo polinomiale.

$$
f \circ g, \qquad f(g(x))
$$
%
è calcolabile in tempo polinomiale.
La composizione di funzioni computabili in tempo polinomiale, è computabile in tempo polinomiale.

\paragraph{Dimostrazione}
Siamo $\mathfrak{M}_f, \, \mathfrak{M}_g$ le macchine di Turing che calcolano rispettivamente $f$ e $g$
\begin{enumerate}
	\item Si copia l'input dal primo al secondo nastro
	\item Si simula $\mathfrak{M}_g$ sul secondo nastro
	\item Se la simulazione termina, si copia l'output sul primo nastro
	\item Si simula $\mathfrak{M}_f$
\end{enumerate}
%
Se anche la seconda simulazione termina, l'output è $f(g(X))$. Osserviamo che se la funzione composta termina, allora $X \in Dom(g)$ e $g(X) \in Dom(f)$, quindi questa macchina di Turing calcola effettivamente la composizione. 

Il tempo polinomiale viene dimostrato nel seguente modo: denotiamo due tempi di computazione sotto forma di polinomi $p_f, \, p_g$. Senza perdere di generalità, possiamo supporre $p_f$ sia un polinomio semplice (crescente), in quanto è sempre possibile maggiorare con un polinomio semplice. Per input $X$:
\begin{enumerate}
	\item La copia iniziale di $X$ richiede tempo lineare in $|X|$
	\item La simulazione di $\mathfrak{M}_g$ richiede al più tempo $p_g(|X|)$
	\item Per il teorema sulla lunghezza dell'output, 
	$$
	|g(X)| \leq |X| + p_g(|X|) + 1,
	$$
	quindi la fase di copia di $g(X)$ dal secondo al primo nastro richiede tempo polinomiale, al più $|X| + p_g(|X|) + 1$. 
	\item La simulazione di $\mathfrak{M}_f$ richiede tempo al più $p_f(|X| + p_g(|X|) + 1)$
\end{enumerate}
%
Il tempo complessivo è quindi limitato dal polinomio

$$
q(|X|) + p_g(|X|) +r(|X| + p_g(|X|) + 1) + p_f(|X| + p_g(|X|) + 1)
$$
%
la computazione della composizione di funzioni computabili in tempo polinomiale, avviene in tempo polinomiale.
Usando il \textbf{Teorema di simulazione di macchine di turing multinastro}, il tempo rimane polinomiale.



